\documentclass[a4paper]{article}
% Español (México) con XeLaTeX: fontspec para fuentes Unicode, polyglossia
% para la división silábica y los nombres del documento en la variante mexicana.
\usepackage{fontspec}
\usepackage{polyglossia}
\setdefaultlanguage[variant=mexican]{spanish}
% Los nombres chinos van como «pinyin (original)»: xeCJK da a los caracteres
% Han su propia fuente; sin ella XeLaTeX los omite con un aviso.
% FandolSong no cubre algunos caracteres de nombres (U+73FA); WenQuanYi Zen Hei sí.
\usepackage[AutoFallBack=true]{xeCJK}
\setCJKmainfont{FandolSong-Regular.otf}
\setCJKfallbackfamilyfont{\CJKrmdefault}{WenQuanYi Zen Hei}
% polyglossia no traduce los nombres de listings.
\AtBeginDocument{\providecommand{\lstlistingname}{}\renewcommand{\lstlistingname}{Listado}%
  \providecommand{\lstlistlistingname}{}\renewcommand{\lstlistlistingname}{Índice de listados}}
\usepackage{amsmath,amssymb}
\usepackage{graphicx}
\usepackage[margin=2.35cm]{geometry}
\usepackage[most]{tcolorbox}
\usepackage{etoolbox}
\usepackage{listings}
\usepackage{xcolor}
\usepackage{hyperref}
\usepackage{booktabs,longtable,tabularx,array}
\usepackage{subcaption}
\usepackage{float}
\usepackage{enumitem}
\usepackage{microtype}
\usepackage{fancyhdr}
\usepackage{needspace}

\definecolor{kbBack}{HTML}{F2F6FA}
\definecolor{kbFrame}{HTML}{9DB4CC}
\definecolor{kbTitle}{HTML}{52708F}
\definecolor{ibBack}{HTML}{FBF5E7}
\definecolor{ibFrame}{HTML}{C9A961}
\definecolor{ibTitle}{HTML}{7D5F1E}
\definecolor{wbBack}{HTML}{FCF2F0}
\definecolor{wbFrame}{HTML}{D6A096}
\definecolor{wbTitle}{HTML}{9E4F43}
\definecolor{dbBack}{HTML}{F4F7F3}
\definecolor{dbFrame}{HTML}{AEC2AC}
\definecolor{dbTitle}{HTML}{5F7F64}

\tcbset{softbox/.style={
  enhanced,breakable,boxrule=0.8pt,arc=2pt,left=3mm,right=3mm,
  top=2mm,bottom=2mm,coltitle=white,fonttitle=\bfseries,
  toptitle=0.6mm,bottomtitle=0.6mm
}}
\newtcolorbox{knowledgebox}[1]{softbox,colback=kbBack,colframe=kbFrame,colbacktitle=kbTitle,title={#1}}
\newtcolorbox{importantbox}[1]{softbox,colback=ibBack,colframe=ibFrame,colbacktitle=ibTitle,title={#1}}
\newtcolorbox{warningbox}[1]{softbox,colback=wbBack,colframe=wbFrame,colbacktitle=wbTitle,title={#1}}
\newtcolorbox{dialoguebox}[1]{softbox,colback=dbBack,colframe=dbFrame,colbacktitle=dbTitle,title={#1}}

\lstset{
  language=Python,basicstyle=\ttfamily\small,
  keywordstyle=\color{kbTitle}\bfseries,stringstyle=\color{wbTitle},
  commentstyle=\color{dbTitle},breaklines=true,frame=single,numbers=left,
  numberstyle=\tiny\color{gray},captionpos=b,columns=fullflexible,
  keepspaces=true,showstringspaces=false,extendedchars=true
}
\BeforeBeginEnvironment{lstlisting}{\Needspace{12\baselineskip}}

\hypersetup{colorlinks=true,linkcolor=kbTitle,urlcolor=kbTitle,citecolor=wbTitle}
\setlist{itemsep=0.25em,topsep=0.4em}
\setlength{\parindent}{2em}
\setlength{\parskip}{0.25em}
\newcolumntype{L}[1]{>{\raggedright\arraybackslash}p{#1}}

\providecommand{\notetitle}{Notas del curso: [tema del curso]}
\providecommand{\noteauthors}{Elaboradas a partir de los materiales públicos de Stanford CS25}
\providecommand{\notedate}{Versión reescrita en 2026}
\providecommand{\videochannel}{Stanford Online}
\providecommand{\videopublishdate}{2022}
\providecommand{\videoduration}{[duración del video]}
\providecommand{\videourl}{}
\providecommand{\videocoverpath}{cover.jpg}
\providecommand{\coursesite}{https://web.stanford.edu/class/cs25/}
\providecommand{\repourl}{https://github.com/hqhq1025/ai-course-notes}
\providecommand{\skillversion}{wdkns/wdkns-skills@39f1a04}

\newcommand{\figsource}[1]{\par\vspace{-0.3em}{\footnotesize\color{gray}\emph{Fuente: #1}}\par}
\newcommand{\teachervoice}[1]{\begin{knowledgebox}{Nota de clase}#1\end{knowledgebox}}
\newcommand{\term}[1]{\textbf{#1}}

\pagestyle{fancy}
\fancyhf{}
\fancyfoot[C]{\thepage}
\fancyfoot[R]{\footnotesize\color{gray}\href{\repourl}{AI Course Notes}}
\renewcommand{\headrulewidth}{0pt}
\renewcommand{\footrulewidth}{0pt}

\newcommand{\makecscover}{
\begin{titlepage}
\centering
\vspace*{0.7cm}
{\Large Stanford CS25 · Transformers United\par}
\vspace{0.8cm}
{\huge\bfseries \notetitle\par}
\vspace{0.6cm}
{\large \noteauthors\par}
\vspace{0.25cm}
{\large \notedate\par}
\vspace{0.9cm}
\ifdefempty{\videocoverpath}{}{
  \includegraphics[width=0.86\textwidth,height=0.43\textheight,keepaspectratio]{\videocoverpath}\par
}
\vfill
\begin{tcolorbox}[width=0.92\textwidth,colback=black!2!white,colframe=black!45,boxrule=0.7pt,arc=2pt]
\textbf{Autor o canal del video}: \videochannel\par
\textbf{Fecha de publicación}: \videopublishdate\par
\textbf{Duración del video}: \videoduration\par
\textbf{Sitio del curso}: \href{\coursesite}{\nolinkurl{\coursesite}}\par
\ifdefempty{\videourl}{}{\textbf{Enlace del video}: \href{\videourl}{\nolinkurl{\videourl}}\par}
\textbf{Normas de elaboración}: \skillversion\ + las normas source-first / teacher-voice / visual-QA de este repositorio
\end{tcolorbox}
\end{titlepage}
\tableofcontents
\newpage
}


\renewcommand{\notetitle}{Lecture 12: Alineación de los modelos de lenguaje con la intención humana}
\renewcommand{\noteauthors}{Elaborado a partir de la conferencia de Jan Leike, los subtítulos oficiales hechos por personas y los artículos originales}
\renewcommand{\notedate}{CS25 V2 · Versión reescrita source-first de 2026}
\renewcommand{\videochannel}{Stanford Online}
\renewcommand{\videopublishdate}{Clase: 2023-01-17; publicación: 2023-05-20}
\renewcommand{\videoduration}{1:06:21}
\renewcommand{\videourl}{https://www.youtube.com/watch?v=DJ1Yy6Aquug}
\renewcommand{\videocoverpath}{cover.jpg}

\newcommand{\lecturefigure}[3]{%
\begin{figure}[H]
  \centering
  \includegraphics[width=0.94\textwidth]{slides-images/#1}
  \caption{#2}
  \figsource{#3}
\end{figure}
}

\begin{document}
\makecscover

\section{Auditoría de fuentes y límites del problema: esta es una clase de alineación de 2023}

Esta sesión se grabó el 17 de enero de 2023, poco después de que ChatGPT se hiciera público. Trata la ruta de RLHF que ya estaba consolidada en ese momento, la scalable oversight que se estaba explorando y las preguntas abiertas que plantearon los estudiantes sobre preferencias, engaño, herramientas, inner alignment e interpretabilidad. La versión anterior trasladaba a la clase hechos posteriores: el relato organizacional de ``Superalignment'', procesos genéricos de puesta en producción y listas de gobernanza de producto. Por eso, esta vez se reescribió por completo con base en el video oficial, 1,264 subtítulos hechos por personas, 17 imágenes didácticas de la clase y seis artículos de primera mano ya publicados en ese momento, en lugar de ampliar la estructura anterior a base de parches.

\lecturefigure{slide-01-title.jpg}{Título de la sesión de Jan Leike: Aligning Language Models with Humans.}{Video oficial de Stanford Online 00:00:23--00:00:49.}

En el título, ``with humans'' es más concreto que ``make models safe'': la sesión pregunta cómo pueden las personas convertir en señal de entrenamiento intenciones que no se pueden escribir por completo como reglas, y cómo mantener útil esa señal mientras la capacidad de los modelos sigue aumentando. El lector debe distinguir tres tipos de evidencia: las slides presentan la línea pública de la exposición, las explicaciones orales del ponente añaden motivaciones y limitaciones, y los artículos de primera mano aportan fórmulas y detalles experimentales. Cuando las tres no coinciden, la nota marca el límite de forma explícita en lugar de hacer pasar conocimiento posterior por palabras textuales de la clase.

\begin{importantbox}{El contrato de evidencia de esta sesión}
Las cifras, el estado de los sistemas y los juicios de investigación que aparecen en la clase tienen como límite el 2023-01-17. El artículo público de InstructGPT puede usarse para aclarar la relación entre SFT, reward model y PPO; los nombres de proyectos, las arquitecturas de despliegue y los cambios de política posteriores a 2023 no pueden presentarse como respuestas que el ponente ya había dado entonces. Durante la última media hora se muestra casi todo el tiempo la última slide, por lo que la teacher voice proviene principalmente de los subtítulos con marca de tiempo, y no de insertar repetidamente la misma imagen del ponente.
\end{importantbox}

\subsection{Team AI: por qué el aumento de capacidades convierte la alineación en un problema estratégico}

El ponente no empieza con funciones de pérdida, sino que plantea una competencia entre ``Team AI y Team Human''. El objetivo de esta analogía no es afirmar que los sistemas actuales ya superan a las personas en todo, sino introducir la dinámica en el problema: distintos AI player van entrando poco a poco, y los player posteriores pueden ser más rápidos, más baratos y más capaces en más tareas. Por ahora, las personas conservan una ventaja importante: decidir qué sistemas entran al mundo real y de qué manera, y entrenar a algunos de ellos para que se pongan del lado humano.

\lecturefigure{slide-02-team-ai-strong-players.jpg}{Team AI se está incorporando a la competencia, y en el futuro podrían aparecer player sumamente fuertes.}{Video oficial de Stanford Online 00:00:49--00:03:18.}

\begin{knowledgebox}{Lectura de la figura: un player fuerte no significa «hoy ya lo puede todo»}
Esta página ve la capacidad de los modelos como una secuencia de player que se incorporan con el tiempo. El ponente subraya a la vez que los sistemas de 2023 todavía tienen limitaciones evidentes; ChatGPT es solo un punto de referencia que muestra cómo la amplitud de conocimientos, la cobertura de idiomas y un costo marginal bajo podrían cambiar la división del trabajo. El ``incredibly strong'' de la figura es una hipótesis de modelado de riesgos orientada al futuro, no una conclusión medida sobre cualquier modelo contemporáneo.
\end{knowledgebox}

Después, la analogía se divide en dos objetivos de acción. Primero, entrenar o elegir AI player dispuestos a ayudar a las personas y a actuar conforme a la intención humana; segundo, diseñar instituciones y reglas para que un entorno competitivo con IA poderosa no lleve a la humanidad a la derrota. El ponente llama alignment al primero; el segundo se acerca más a governance. Ambos se influyen mutuamente, pero no son el mismo problema técnico.

\lecturefigure{slide-03-two-objectives.jpg}{Los dos objetivos de Team Human: reclutar AI player y establecer reglas con las que no se pierda la competencia.}{Video oficial de Stanford Online 00:03:18--00:04:16.}

\begin{warningbox}{No mezclar alignment y governance en un término que lo abarque todo}
Alignment estudia señales de entrenamiento, objetivos de comportamiento, evaluación y generalización; governance se ocupa de los permisos de despliegue, las reglas de competencia, la responsabilidad y la coordinación institucional. Esta sesión solo desarrolla lo primero. Introducir en lo que sigue procesos organizacionales, propuestas regulatorias o la competencia de mercado no solo se aparta de la clase, sino que también oculta el verdadero cuello de botella técnico de la sesión: cuando las personas no entienden las tareas que realiza el modelo, cómo seguir dando retroalimentación confiable.
\end{warningbox}

\subsection{Resumen de la sección}

Esta sección establece dos límites. En el tiempo, se trata de una clase de principios de 2023; en el problema, se centra en cómo lograr que los modelos sigan la human intent, sin intentar resolver toda la AI governance. La analogía de Team AI aporta la motivación estratégica; lo siguiente es descomponer ``ponerse del lado humano'' para que deje de ser un lema y se convierta en objetivos entrenables.
\section{Human intent: por qué ejecutar al pie de la letra sigue sin ser suficiente}

La sección anterior explicó por qué se necesita alignment; esta sección se pregunta cuál es exactamente el objetivo. Un modelo puede ejecutar una orden palabra por palabra y, aun así, contradecir la intención real al inventar hechos, aprovechar ambigüedades, ignorar los límites de seguridad o complacer la formulación inmediata del usuario. Por eso, el núcleo de alignment no es la ``obedience'', sino combinar la solicitud explícita, las restricciones implícitas y el manejo de la incertidumbre.

\subsection{Del eslogan a la definición operativa}

La formulación práctica que da el ponente es: construir sistemas que sigan la human intent y las human preferences. Esta definición evita a propósito expresar la intención como una regla cerrada, porque los usuarios reales suelen aportar solo información parcial. El modelo tiene que inferir el resultado que el usuario realmente quiere obtener a partir de la semántica de la tarea, el historial de la conversación, el sentido común general y las restricciones de seguridad. A continuación, primero se usa la diapositiva de título para fijar este problema operativo y después se descompone la intent en partes explícitas, implícitas y pendientes de aclarar; al leer la figura, conviene notar en especial que ``follow'' apunta a resultados y normas, no solo a repetir el texto superficial de la orden.

\lecturefigure{slide-04-follow-human-intent.jpg}{El problema operativo de alignment: ¿cómo construir sistemas de IA que sigan la intención humana?}{Video oficial de Stanford Online, 00:04:16--00:05:31.}

\begin{importantbox}{Descomposición mínima de la human intent}
El objetivo de una interacción puede escribirse como
\[
I(x,c)=I_{\mathrm{explicit}}(x,c)+I_{\mathrm{implicit}}(x,c)+I_{\mathrm{uncertain}}(x,c),
\]
donde $x$ representa la solicitud actual del usuario y $c$ el contexto de la conversación y del entorno; $I_{\mathrm{explicit}}$ es la tarea literal, $I_{\mathrm{implicit}}$ son las normas de veracidad, seguridad y cooperación que no se dicen pero se esperan razonablemente, e $I_{\mathrm{uncertain}}$ representa la parte que el modelo todavía no puede juzgar con fiabilidad y que debe manejar mediante una aclaración o una acción conservadora. Esta suma es una descomposición didáctica; no implica que los tres términos puedan medirse de forma exacta e independiente.
\end{importantbox}

\subsection{Explicit intent e implicit intent}

La intención explícita incluye instrucciones directas como ``resume este artículo'' o ``actúa como mi asistente''; la intención implícita incluye no inventar, no tergiversar deliberadamente, rechazar solicitudes dañinas y preguntar cuando hay una ambigüedad en un punto crítico. Esta última es difícil porque, en la cooperación cotidiana, las personas dependen de una gran cantidad de normas implícitas y no reescriben un manual de normas completo en cada solicitud.

\lecturefigure{slide-05-explicit-implicit-intent.jpg}{Intención explícita e intención implícita: seguir instrucciones es solo una parte del objetivo de alineación.}{Video oficial de Stanford Online, 00:05:31--00:06:35.}

\begin{knowledgebox}{Lectura de la figura: implicit no equivale a «que el modelo adivine libremente»}
La columna izquierda puede observarse directamente en la instrucción o en el rol; la columna derecha debe inferirse de las normas de la interacción. Un sistema razonable debe tratar las restricciones implícitas como hipótesis con un grado de incertidumbre: ante una ambigüedad de bajo riesgo puede seguir con la tarea; ante una ambigüedad de alto riesgo debe pedir una aclaración. Por lo tanto, ``ask follow-up questions'' no es un adorno de la conversación, sino una acción de medición activa que convierte una intención no observable en retroalimentación observable.
\end{knowledgebox}

\begin{warningbox}{La puerta de entrada al specification gaming}
Si el entrenamiento solo recompensa el cumplimiento literal, el sistema aprenderá a aprovechar las zonas que la métrica de evaluación no cubre. Por ejemplo, un resumen que parece fluido pero inventa fuentes, un código que pasa las pruebas públicas pero oculta una puerta trasera, o un asistente que da respuestas que al usuario le gusta oír pero que no son ciertas. El problema no es que el modelo ``no haya ejecutado la orden'', sino que la orden y la evaluación solo describen una parte de la intención real.
\end{warningbox}

\subsection{Resumen de la sección}

La human intent comprende, como mínimo, la tarea explícita, las normas implícitas y el manejo de la incertidumbre. Los métodos de alineación deben aprender, a partir de una retroalimentación limitada, estas preferencias que nunca se escriben por completo, y a la vez evitar confundir los hábitos de un solo evaluator con el objetivo único de toda la humanidad. RLHF es precisamente el mecanismo de conversión más representativo de este tipo en los modelos actuales.
\section{RLHF: convertir comparaciones humanas en actualizaciones de la política}

Esta sección parte del diagrama de dos pasos de la slide y reconstruye la cadena didáctica completa SFT--RM--RL. La idea central es la siguiente: las personas no necesitan asignar una puntuación absoluta a cada salida, sino solo comparar unos pocos candidatos; el reward model aprende el patrón de ese ordenamiento y, después, proporciona una señal aproximada y optimizable para una gran cantidad de muestras sin anotación humana.

\subsection{Primer paso: aprender un reward model a partir de pairwise comparison}

Dado un prompt $x$, la política produce dos respuestas, $y_w$ y $y_l$, y el anotador elige como winner la que mejor se ajusta a la intención. El reward model $r_\phi(x,y)$ no genera texto directamente, sino que asigna un escalar a cada par prompt--response. La forma de Bradley--Terry, de uso común, convierte la diferencia de recompensas en una probabilidad de preferencia:

\[
P_\phi(y_w \succ y_l\mid x)
=\sigma\!\left(r_\phi(x,y_w)-r_\phi(x,y_l)\right),
\]

donde $\sigma(z)=1/(1+e^{-z})$ es la función logística; $r_\phi$ es el reward model con parámetros $\phi$; el símbolo $y_w \succ y_l$ indica que el anotador prefiere $y_w$. La log-verosimilitud negativa correspondiente es

\[
\mathcal{L}_{\mathrm{RM}}(\phi)
=-\mathbb{E}_{(x,y_w,y_l)\sim D}
\log \sigma\!\left(r_\phi(x,y_w)-r_\phi(x,y_l)\right).
\]

Aquí, $D$ denota el conjunto de datos de comparaciones humanas, y la esperanza promedia sobre los prompts y los pares de respuestas. La pérdida solo restringe el orden relativo; por ello, el cero absoluto de la recompensa carece de significado propio.

\lecturefigure{slide-06-rlhf-reward-model.jpg}{Primer paso de RLHF: entrenar un reward model a partir de comparaciones humanas.}{Video oficial de Stanford Online, 00:06:35--00:07:55.}

\begin{knowledgebox}{Lectura de la figura: qué problema resuelven los datos de comparación}
En la figura, las personas evalúan varios response para un mismo prompt y proporcionan un orden relativo; el reward model aprende a predecir ese orden. Un juicio relativo suele ser más natural que escribir una recompensa precisa, y además reparte el costo de una sola comparación humana entre muchas policy samples posteriores. No elimina la subjetividad: distintos labeler discrepan, y lo que el modelo aprende al final es la preferencia agregada de la distribución de los datos, no una lectura directa de una verdad objetiva.
\end{knowledgebox}

\teachervoice{El ponente afirma explícitamente que los labeler a menudo no coinciden y que el modelo agrega esas diferencias durante el entrenamiento. El problema de ingeniería aquí no se reduce al ``ruido en las etiquetas'': también incluye a quién se contrata para anotar, qué lineamientos se le dan, qué grupos quedan fuera de los datos y si las preferencias en conflicto deberían promediarse.}

\subsection{Segundo paso: optimizar la policy y limitar a la vez su desviación}

Una vez obtenido el reward model, se pueden muestrear respuestas de la política y usar aprendizaje por refuerzo para aumentar la recompensa predicha. Si solo se maximiza $r_\phi$, la policy puede aprovechar los huecos del reward model; por eso InstructGPT emplea una restricción de KL respecto de una política de referencia y mezcla la pretraining loss para reducir la regresión de capacidades. Un objetivo con fines didácticos es

\[
J(\theta)=
\mathbb{E}_{x\sim D_x,\,y\sim\pi_\theta(\cdot\mid x)}
\left[r_\phi(x,y)-\beta\,
D_{\mathrm{KL}}\!\left(\pi_\theta(\cdot\mid x)\,\|\,\pi_{\mathrm{ref}}(\cdot\mid x)\right)\right]
+\gamma\,\mathbb{E}_{z\sim D_{\mathrm{ptx}}}\log \pi_\theta(z).
\]

Aquí, $\pi_\theta$ es la policy por optimizar, $\pi_{\mathrm{ref}}$ es la SFT/reference policy, $\beta$ controla la intensidad de la penalización por alejarse del modelo de referencia, $D_{\mathrm{KL}}$ mide la diferencia entre las dos distribuciones de salida, $D_x$ es la distribución de prompts, $D_{\mathrm{ptx}}$ es el texto de preentrenamiento y $\gamma$ controla el peso que se da a conservar la capacidad de modelado del lenguaje.

\lecturefigure{slide-07-rlhf-full-loop.jpg}{Narrativa completa de la slide: el reward model proporciona retroalimentación y el modelo de lenguaje preentrenado obtiene una nueva policy mediante RL.}{Video oficial de Stanford Online, 00:07:55--00:09:24.}

\begin{importantbox}{El diagrama de dos pasos de la clase y las tres etapas del artículo no se contradicen}
Para destacar el reward modeling, la slide condensa el proceso en ``entrenar el RM con comparaciones'' y ``optimizar la policy con RL''. El ponente advierte de palabra que se omitió un paso. Según el artículo de InstructGPT, el proceso habitual es: primero, supervised fine-tuning (SFT) con demonstrations; después, entrenar el reward model (RM) con comparisons, y, por último, optimizar la policy con un método de RL como PPO. Que estas notas añadan el SFT es una aclaración de la fuente, no una pretensión de que en la slide aparezca un tercer recuadro.
\end{importantbox}

\begin{lstlisting}[caption={Pseudocódigo mínimo de las tres etapas de RLHF}]
policy = supervised_finetune(base_model, demonstrations)
reward_model = train_pairwise_reward_model(policy, comparisons)

for batch in prompt_batches:
    responses = policy.sample(batch)
    rewards = reward_model.score(batch, responses)
    policy = ppo_update(
        policy,
        rewards=rewards,
        reference=supervised_policy,
        kl_penalty=beta,
        pretraining_batch=next(pretraining_data),
    )
\end{lstlisting}

\subsection{Por qué no limitarse a imitation learning}

El SFT/behavioral cloning aprende ``cómo respondería una persona''; la preference optimization aprende ``qué resultado prefieren las personas''. El ponente señala una limitación en dos sentidos: en algunas tareas las personas superan al modelo, y este quizá no logre imitarlas a la perfección; en otras, el modelo ya tiene ventajas distintas, y obligarlo a imitar la secuencia de acciones humana más bien lo limita. El RL permite que el modelo encuentre su propia manera de resolver la tarea, siempre que el resultado supere la evaluación por preferencias.

\begin{knowledgebox}{Términos clave: SFT, RM, PPO, PPO-ptx}
\begin{tabularx}{\textwidth}{L{2.2cm} >{\raggedright\arraybackslash}X >{\raggedright\arraybackslash}X}
\toprule
Término & Problema que resuelve & Lugar en esta clase \\
\midrule
SFT & Construye, con demonstrations de alta calidad, una policy inicial con la que se puede interactuar & Proporciona la política de referencia y la capacidad básica de seguir instrucciones \\
RM & Condensa las preferencias por pares en una función escalar que se puede invocar repetidamente & Permite que un número limitado de comparaciones humanas supervise una gran cantidad de muestras \\
PPO & Aumenta la recompensa predicha limitando la magnitud de cada actualización & Era un optimizador conocido y eficaz en ese momento, no uno cuya optimalidad se haya demostrado \\
PPO-ptx & Mezcla el objetivo de preentrenamiento durante el RL & Mitiga la regresión de capacidades que causa perseguir solo la recompensa por preferencias \\
\bottomrule
\end{tabularx}
\end{knowledgebox}

\teachervoice{Un estudiante pregunta por qué se usa PPO. La respuesta del ponente es muy de ingeniería: el equipo lo conocía bien y ``works pretty well''. Esto muestra que la elección de un algoritmo depende de la trayectoria histórica y de la madurez de las herramientas; de ``fue adoptado'' no se puede deducir ``es el mejor en teoría''.}

\subsection{Resumen de la sección}

RLHF entrena un reward model con comparaciones humanas y después optimiza la policy; en la práctica completa incluye además SFT, KL regularization y la conservación de datos de preentrenamiento. Convierte el costo de la retroalimentación, que antes era una calificación humana por cada acción, en un modelo reutilizable, pero también introduce en el objetivo de entrenamiento las preferencias de los labeler, el reward hacking y la proxy misspecification.
\section{InstructGPT y el primer ChatGPT: efectos, costos y límites}

Para saber si RLHF vale la pena no basta con mirar la función de pérdida: también hay que ver qué prefieren los usuarios, cuántos recursos se invierten y si las capacidades previas retroceden. Esta sección examina, en orden, tres evidencias presentadas en clase: la curva de preferencias, el costo de entrenamiento y el balance de la experiencia con ChatGPT a principios de 2023.

\subsection{100 veces más pequeño, ¿por qué sigue siendo preferido?}

El artículo de InstructGPT somete modelos base GPT-3, SFT, PPO y PPO-ptx de distintos tamaños a la misma evaluación de preferencias humanas. El eje horizontal es el tamaño del modelo en parámetros y el vertical, la frecuencia con que se elige frente a cierta línea base. El resultado más llamativo es que, en esa distribución de customer-prompt, la variante de InstructGPT de 1.3B puede preferirse sobre el GPT-3 base de 175B.

\lecturefigure{slide-08-instructgpt-preferences.jpg}{Curva de preferencias de InstructGPT: el alignment fine-tuning puede cambiar de forma notable la utilidad que perciben los usuarios.}{Video oficial de Stanford Online 00:09:24--00:15:18.}

\begin{knowledgebox}{Lectura de la figura: primero el eje y, después el tamaño del modelo}
El eje y no mide intelligence general ni el promedio de todos los benchmarks, sino la preference win rate humana sobre una distribución específica de prompts. La figura respalda que «un aligned model más pequeño puede ajustarse mejor a la intención del usuario»; no respalda que «el modelo de 1.3B supera al de 175B en conocimiento, razonamiento y todas las tareas». La curva de PPO-ptx recuerda además que la ganancia en preferencia y la conservación de capacidades deben medirse al mismo tiempo.
\end{knowledgebox}

Si el experimento de preferencias se formula como una variable aleatoria binomial, la tasa de victorias estimada del modelo $A$ frente al modelo $B$ es

\[
\widehat{p}_{A>B}=\frac{1}{N}\sum_{i=1}^{N}\mathbb{1}[A\text{ se elige en la comparación }i],
\]

donde $N$ es el número de comparaciones válidas y $\mathbb{1}[\cdot]$ es la función indicadora. Este estadístico depende de la distribución de prompts, de los labelers, del orden de presentación y de los criterios de evaluación; con otro conjunto de tareas u otra población, $\widehat{p}_{A>B}$ puede cambiar.

\begin{warningbox}{Preference no es sinónimo de capability}
La preferencia puede premiar la claridad, la cooperación, el rechazo de solicitudes indebidas y el apego al formato, pero también la seguridad aparente, la complacencia o un estilo familiar. Por eso, la preference win debe considerarse junto con la veracidad, la calibración, la robustez, la exactitud en tareas especializadas y las evaluaciones de seguridad, y no puede cargar por sí sola con todo el significado de «el modelo es mejor».
\end{warningbox}

\subsection{El balance de recursos del alignment fine-tuning}

La slide de costos coloca en un mismo eje vertical el compute de GPT-3 pretraining, de InstructGPT SFT y de InstructGPT RL, y muestra aparte, a la derecha, la retroalimentación humana. Lo que expone no es que «el alignment sea gratis», sino que las distintas dimensiones de recursos son asimétricas: frente al preentrenamiento, el cómputo de SFT/RL es muy pequeño; pero reunir demonstrations y comparisons de alta calidad exige mucho tiempo humano, diseño de criterios y control de calidad.

\lecturefigure{slide-09-instructgpt-training-costs.jpg}{Costo de entrenamiento de InstructGPT: el alignment compute es muy pequeño, pero la retroalimentación humana aún requiere unas 20,000 horas.}{Video oficial de Stanford Online 00:15:18--00:19:05.}

\begin{knowledgebox}{Lectura de la figura: no sumes directamente PFLOP/s-days y dólares}
La gráfica de barras de la izquierda compara cómputo; el texto de la derecha da la estimación del ponente para la retroalimentación humana: unos 500,000 dólares y unas 20,000 horas. Son recursos distintos: el compute puede paralelizarse con hardware, mientras que el juicio humano está sujeto a restricciones de reclutamiento, capacitación, consistencia y representatividad. El costo real del sistema incluye también la gobernanza de datos, el diseño de evaluaciones, la iteración de ingeniería y los experimentos fallidos; la slide no afirma cubrir el costo de todo el ciclo de vida.
\end{knowledgebox}

\begin{table}[H]
\centering
\caption{Tres tipos de recursos que no pueden reducirse a una sola cifra de precio}
\begin{tabularx}{\textwidth}{L{2.9cm} X X}
\toprule
Recurso & Función principal & Riesgo principal \\
\midrule
Pretraining compute & Adquirir capacidades amplias de lenguaje y de tareas & Costo alto, límites de datos opacos, fecha de corte del conocimiento \\
Alignment compute & Hacer las capacidades existentes más controlables y más acordes con los objetivos de interacción & Reward hacking, retroceso de capacidades, sobreajuste a las preferencias \\
Human feedback & Aportar demonstrations, comparisons y juicios sobre errores & Representatividad insuficiente, deriva de criterios, fatiga y problemas de consistencia \\
\bottomrule
\end{tabularx}
\end{table}

\subsection{ChatGPT: la interfaz de diálogo mejora la usabilidad, pero no elimina el desajuste}

El ponente presenta ChatGPT como el siguiente paso de InstructGPT: el diálogo se vuelve una interfaz general, el usuario puede repreguntar, corregir y pedir que se reescriba una respuesta; también mejora el rechazo de tareas dañinas. Pero eso no significa que «la alineación esté terminada». A principios de 2023, las limitaciones más notorias seguían siendo la hallucination y la prompt sensitivity: el modelo inventa hechos, y distintas formulaciones pueden provocar respuestas de distinta calidad o comportamientos de seguridad distintos.

\lecturefigure{slide-10-chatgpt-lessons.jpg}{Balance de ChatGPT presentado en clase a principios de 2023: dialog, tareas complejas, hallucination y prompt sensitivity.}{Video oficial de Stanford Online 00:19:05--00:20:19.}

\begin{warningbox}{Límite temporal: no expliques esta slide con el ChatGPT posterior}
Esta página describe el sistema público del 2023-01-17. En la sesión de Q\&A, el ponente dice explícitamente que la cantidad exacta de datos y los pasos de entrenamiento de ChatGPT no se han publicado, y que solo puede afirmarse que el proceso es, en general, similar al de InstructGPT; las cifras que recuerda de comparisons, demonstrations y las 20,000 horas también son aproximadas. Los apuntes no deben atribuir retroactivamente a ese momento, como especificaciones conocidas, las capacidades posteriores del producto, el uso de herramientas ni los cambios de política.
\end{warningbox}

\begin{importantbox}{Por qué la prompt sensitivity también es un problema de alignment}
Si dos prompts semánticamente equivalentes provocan, solo por diferencias de redacción, niveles completamente distintos de veracidad, rechazo o calidad de razonamiento, lo que el modelo aprendió sigue siendo un patrón superficial y no una intención de tarea estable. Un sistema ideal debería conservar una personalización razonable y, a la vez, mantener sin cambios su comportamiento esencial ante transformaciones que no alteran el significado, además de preguntar activamente cuando haya ambigüedad.
\end{importantbox}

\subsection{Resumen de la sección}

InstructGPT demuestra que el alignment fine-tuning puede mejorar de forma notable la preferencia humana en una distribución específica, con un costo de cómputo muy inferior al del preentrenamiento; pero el trabajo de retroalimentación, la representatividad y el riesgo de proxy siguen siendo costosos. El primer ChatGPT demostró además que la interfaz de diálogo por sí misma mejora la usabilidad, y a la vez dejó ver que la hallucination y la prompt sensitivity siguen sin resolverse.
\section{Evaluation is easier than generation: la palanca de la supervisión escalable}

Los métodos anteriores dependen de comparaciones humanas. Para extenderlos a modelos más capaces hay que responder una pregunta central: ¿pueden las personas evaluar tareas que ellas mismas no sabrían realizar? El punto de partida del ponente es una asimetría frecuente, aunque no absoluta: generar un resultado de alta calidad puede ser difícil, mientras que, una vez vistos los candidatos, identificar cuál es mejor suele ser más sencillo.

\subsection{De la proposición a la intuición}

La clase se detiene primero en una diapositiva de título aparte porque esta frase sostiene toda la lógica de la segunda mitad. Si la evaluation fuera tan difícil como la generation, las personas perderían todo punto de entrada a la supervisión en cuanto la capacidad del modelo las superara; si evaluar es más fácil, las personas pueden apoyarse en candidatos, evidencia y críticas para extender su capacidad de supervisión a tareas que no podrían completar por sí solas. Esta sección formula primero el principio como una ventaja de costo condicionada y después, mediante cuatro analogías, examina cuándo se cumple y cuándo deja de cumplirse porque la evidencia está oculta o porque el evaluator es manipulado.

\lecturefigure{slide-11-evaluation-title.jpg}{Principio central de la segunda mitad de la clase: evaluation is easier than generation.}{Video oficial de Stanford Online, 00:20:19--00:20:32.}

\begin{importantbox}{Formulación condicionada de la ventaja de evaluar}
Una formulación más precisa es: dados los resultados candidatos, la especificación de la tarea y evidencia suficientemente verificable, el costo de comparar o de revisar partes concretas puede ser menor que el de generar desde cero. La ventaja proviene de que los candidatos reducen el espacio de búsqueda, no de que el evaluator posea automáticamente la verdad. Si el error está oculto en un estado no visible, si la evidencia no puede obtenerse, si se puede persuadir al evaluador o si el propio objetivo de la tarea es ambiguo, la ventaja de evaluar se reduce e incluso desaparece.
\end{importantbox}

\subsection{Cómo deben leerse las cuatro analogías}

La slide enumera P/NP, ver una competencia deportiva, comprar un producto y la revisión académica por pares. Comparten una estructura: quien ejecuta necesita encontrar una solución dentro de un espacio de acciones enorme, mientras que quien evalúa, al ver unos pocos candidatos, solo necesita comparar resultados o señalar problemas concretos. Sin embargo, estos ejemplos no demuestran que «todas las tareas reales puedan verificarse con eficiencia»; solo explican por qué la preference comparison puede ser más escalable que la demonstration.

\lecturefigure{slide-12-evaluation-easier-examples.jpg}{Analogías de que evaluar es más fácil que generar: complejidad computacional, deportes, decisiones de consumo y revisión de artículos.}{Video oficial de Stanford Online, 00:20:32--00:22:25.}

\begin{knowledgebox}{Lectura de la figura: reducir las analogías a una estructura común}
P/NP subraya la diferencia entre encontrar una solución y verificar un certificado; los deportes subrayan que el público puede juzgar quién gana sin alcanzar el nivel profesional; las decisiones de consumo subrayan que comparar la experiencia de uso no exige fabricar el producto; la revisión por pares subraya que señalar defectos no equivale a escribir un artículo mejor. Lo común es que el evaluator dispone de candidatos o de evidencia. La diferencia es que la evaluación en la realidad suele tener objetivos subjetivos, variables ocultas y evidencia incompleta, por lo que no se pueden aplicar directamente las conclusiones de la teoría de la complejidad.
\end{knowledgebox}

Si el costo de generación se denota $C_{\mathrm{gen}}(T)$ y el costo de evaluación con información auxiliar $a$ se denota $C_{\mathrm{eval}}(T\mid a)$, la scalable oversight busca construir $a$ de modo que

\[
C_{\mathrm{eval}}(T\mid a) \ll C_{\mathrm{gen}}(T),
\]

donde $T$ representa la tarea y $a$ puede ser una respuesta candidata, una cita, una prueba, una critique, un diálogo o el resultado de una herramienta. El símbolo $\ll$ es un objetivo de diseño, no un teorema que se cumpla para todas las tareas.

\begin{warningbox}{El evaluador también es objeto de optimización}
En cuanto el modelo sabe qué obtiene una calificación alta, ya no optimiza solo la tarea, sino también el proceso de observación del evaluator. La expresión fluida, la evidencia seleccionada a conveniencia, la complacencia con sesgos e incluso el ocultamiento de defectos pueden elevar la calificación. Que evaluar sea más fácil solo se convierte en una señal de entrenamiento confiable cuando el propio proceso de evaluación es suficientemente robusto frente a estas estrategias.
\end{warningbox}

\subsection{Resumen de la sección}

La ventaja de evaluar ofrece un punto de entrada a la scalable oversight: las personas no necesitan generar por sí solas resultados de nivel experto y aun así pueden comparar candidatos o revisar evidencia concreta. Pero esta ventaja depende de evidencia visible, especificaciones claras y un proceso de evaluación resistente a la manipulación; no es una garantía de que «las personas siempre podrán supervisar a una IA más capaz».
\section{Scalable oversight: cuando el modelo supera el techo de la evaluación humana}

El RLHF convencional supone que los humanos pueden comparar de manera confiable las policy outputs. A medida que la capacidad de los modelos sigue creciendo, este supuesto es el primero en fallar: el modelo puede completar tareas más complejas, mientras que el conocimiento especializado y la atención del evaluator humano no crecen automáticamente con el AI progress. Esta sección describe primero esa falla y después explica cómo la AI assistance intenta elevar la curva de evaluación.

\subsection{Human evaluation ceiling y deception incentive}

En la figura, el eje horizontal es el AI progress y el eje vertical es la dificultad de la tarea. La curva naranja de capacidad asciende, mientras que la unaided human evaluation, en verde, es aproximadamente horizontal. Cuando la línea naranja rebasa la verde, los humanos todavía pueden expresar una preferencia, pero no pueden saber si la salida con puntuación alta es realmente correcta; la reward optimization convierte ese punto ciego en un espacio explotable.

\lecturefigure{slide-13-scaling-human-supervision.jpg}{Figura central de scalable oversight: la capacidad de la IA crece, pero la capacidad de evaluación independiente de los humanos tiene un techo.}{Video oficial de Stanford Online 00:22:25--00:25:30.}

\begin{knowledgebox}{Lectura de la figura: qué representa cada una de las tres curvas}
La curva inferior es la dificultad de las tareas que la IA puede completar; la línea horizontal es el límite superior de lo que los humanos pueden evaluar de forma independiente y confiable; la nueva curva superior es la capacidad de evaluación de los humanos en colaboración con la IA. Lo importante no es el punto de intersección exacto, sino el orden: primero aparecen tareas que el modelo puede realizar pero que los humanos no comprenden, y después surge la optimización sistemática contra una señal de evaluación débil. La figura no demuestra que la deception ocurra necesariamente, pero explica por qué la deception o la sycophancy podrían convertirse en estrategias de alta recompensa.
\end{knowledgebox}

Escribamos la calidad real de la tarea como $q(x,y)$, la puntuación visible para el humano como $h(x,y)$, y el reward model aprende $r_\phi(x,y)\approx h(x,y)$. Si en la región que los humanos no comprenden $h$ se desacopla de $q$, entonces

\[
\arg\max_y r_\phi(x,y) \neq \arg\max_y q(x,y).
\]

Aquí, el lado izquierdo es el proxy que la policy optimiza en realidad, y el lado derecho es la calidad de la tarea que de verdad les importa a los humanos. Cuanto mayor es la capacidad, más probable es que el modelo encuentre salidas complejas que separen ambos.

\begin{warningbox}{«Engañar al evaluator» no requiere que el modelo tenga malicia de tipo humano}
Basta con que el entrenamiento seleccione de forma continua las salidas con puntuación alta para que se refuerce cualquier patrón que eleve la puntuación de manera estable sin mejorar la calidad real. Puede tratarse de complacencia, exceso de confianza, ocultamiento de la incertidumbre, presentación selectiva de evidencia o un engaño estratégico genuino. El riesgo proviene de la estructura de la optimización y la medición, no de suponer de antemano que el modelo posee algún tipo de personalidad.
\end{warningbox}

\subsection{AI-assisted evaluation: descomponer la tarea global en problemas locales verificables}

El ponente pone como ejemplo la auditoría de código: un humano no puede encontrar en un tiempo limitado todos los bugs o Trojans de una codebase grande, pero si el modelo señala la ubicación concreta, las condiciones de falla y las pruebas, el humano puede verificar esa claim local. El papel de la AI assistance no es emitir el juicio de valor final en lugar del humano, sino encargarse de la lectura, la búsqueda, el cálculo, el fact checking y la generación de candidatos, para transformar un objeto difícil de evaluar en un conjunto de unidades de evidencia más pequeñas.

\lecturefigure{slide-14-critiques-and-dialog.jpg}{Dos tipos de AI assistance: generar una critique o consultar evidencia activamente mediante dialogue.}{Video oficial de Stanford Online 00:25:30--00:26:38.}

\begin{knowledgebox}{Lectura de la figura: critique y dialogue ofrecen interfaces de información distintas}
La critique indica de una sola vez ``dónde podría haber un error'', lo que la hace adecuada para exponer puntos de revisión con alta exhaustividad; el dialogue permite que el evaluator pregunte por fuentes, contraejemplos, procedimientos de cálculo y explicaciones alternativas, por lo que es más adecuado para reducir la incertidumbre paso a paso. Ambos dependen de la honestidad y la cobertura del assistant; que su lenguaje suene natural no basta para tratarlos como evidencia confiable.
\end{knowledgebox}

\begin{table}[H]
\centering
\caption{Acciones comunes de la AI-assisted evaluation y responsabilidades restantes}
\begin{tabularx}{\textwidth}{L{2.6cm} X X}
\toprule
Acción & Trabajo cognitivo que puede asumir la IA & Lo que el humano aún debe juzgar \\
\midrule
Critique & Enumerar defectos, contraejemplos, omisiones y contradicciones & Qué problemas son importantes y si la critique es válida \\
Dialogue & Responder preguntas de seguimiento, comparar alternativas, explicar claims locales & Cuándo dejar de preguntar y si la evidencia es suficiente \\
Fact checking & Buscar fuentes, localizar citas, hacer verificaciones cruzadas & Credibilidad de las fuentes y si el contexto corresponde \\
Code/tool use & Ejecutar pruebas, calcular, hacer análisis estático, navegar & Límites de permisos, relación entre los resultados de las herramientas y el objetivo real \\
\bottomrule
\end{tabularx}
\end{table}

\subsection{Resumen de la sección}

Cuando la capacidad del modelo supera la unaided human evaluation, el proxy del RLHF convencional puede desacoplarse de la calidad real. La scalable oversight intenta ayudar a los humanos a evaluar aprovechando las capacidades de lectura, búsqueda, cálculo y crítica de una IA del mismo tipo, pero el juicio final de preferencia, la confianza en la evidencia y los límites de permisos siguen siendo responsabilidad humana.
\section{Cómo medir si la oversight realmente mejoró}

«Pedirle a la IA que ayude a evaluar» es fácil de demostrar, pero difícil de validar científicamente. Las tareas realmente difíciles a menudo no tienen una ground truth disponible; si el evaluator y el assistant llegan a la misma conclusión, no sabemos si descubrieron juntos la verdad o si los convenció a ambos el mismo sesgo. Esta sección presenta dos herramientas de diagnóstico: las targeted perturbations y el discriminator--critique gap.

\subsection{Targeted perturbations: crear defectos conocidos de forma artificial}

El experimento toma primero una response y la usa provisionalmente como reference; después, una persona le inserta una subtle but important flaw y se obtiene una perturbed response. Como los investigadores saben cuál de las dos fue modificada, pueden mostrar al azar una de ellas y medir la diferencia de exactitud con la que un unaided human y un assisted human identifican el defecto.

\lecturefigure{slide-15-measuring-progress.jpg}{Medición de la AI assistance: a partir de la respuesta original se construyen pares de muestras con una subtle flaw conocida.}{Video oficial de Stanford Online 00:26:38--00:28:57.}

\begin{knowledgebox}{Lectura de la figura: qué es exactamente la ground truth}
En este procedimiento, la ground truth es «cuál de las dos fue perturbada artificialmente», no que la respuesta original ya esté demostrada como absolutamente correcta. Este diseño permite que el experimento mida si la assistance facilita detectar el defecto insertado, al tiempo que conserva la complejidad de lenguaje natural de la tarea; no puede abarcar todos los patrones de error que surgen de forma natural y que los investigadores no anticiparon.
\end{knowledgebox}

Sea $z_i=1$ si la $i$-ésima muestra se clasificó correctamente como flawed/correct, y sea $a\in\{0,1\}$ la condición de asistencia; entonces

\[
\widehat{\mathrm{Acc}}_a=\frac{1}{N_a}\sum_{i:a_i=a} z_i,
\qquad
\Delta_{\mathrm{assist}}=\widehat{\mathrm{Acc}}_1-\widehat{\mathrm{Acc}}_0.
\]

Aquí, $N_a$ es el número de muestras de la condición correspondiente, $\widehat{\mathrm{Acc}}_1$ es la tasa de aciertos con AI assistance, $\widehat{\mathrm{Acc}}_0$ es la tasa de aciertos sin assistance y $\Delta_{\mathrm{assist}}$ mide la ganancia neta. También deben reportarse los falsos positivos, la gravedad del defecto y si el evaluator habría notado el problema por su cuenta.

\begin{lstlisting}[caption={Pseudocódigo mínimo de la evaluación con targeted perturbations}]
for prompt, reference in task_set:
    flawed = human_insert_subtle_flaw(reference)
    shown, label = random_choice([
        (reference, "reference"),
        (flawed, "flawed"),
    ])

    unaided = human_judge(shown)
    critique = assistant_critique(prompt, shown)
    assisted = human_judge(shown, critique=critique)

    record(label, unaided, assisted, critique)
\end{lstlisting}

\subsection{Critique experiment: prometedor, pero lejos de cerrar el ciclo}

Esta subsección usa el experimento de topic-based summarization para poner a prueba el método de evaluación de la subsección anterior: las model critiques ayudan a las personas a encontrar alrededor de 50\% más flaws. Este resultado muestra que el modelo puede convertir parte de su conocimiento de la tarea en pistas de revisión que el evaluator puede usar; sin embargo, el ponente plantea de inmediato tres limitaciones: la tarea todavía no es lo bastante difícil, el efecto real tendría que ser mayor para resultar tranquilizador y muchas critiques son basura o nitpicking. Por ello, el orden de lectura de la figura debe ser primero confirmar la diferencia relativa con y sin assistance, y luego preguntar por la tasa absoluta de defectos no detectados, la gravedad de los defectos y el alcance de la extrapolación a otras tareas, en lugar de quedarse solo con el headline percentage. Solo leyendo el tamaño del efecto junto con las muestras fallidas se puede juzgar si la assistance aporta una comprensión real o si solo aumenta la cantidad de texto que el evaluator tiene que descartar.

\lecturefigure{slide-16-summarization-critiques-results.jpg}{Resultados presentados en clase de las summarization critiques: la assistance aumentó la cantidad de defectos que encuentran las personas.}{Video oficial de Stanford Online 00:28:57--00:29:45.}

\begin{warningbox}{Tres limitaciones que no deben omitirse al leer la figura}
Primero, «50\% more flaws» es un incremento relativo, no equivale a encontrar todos los defectos; segundo, la summarization es más fácil de verificar que las futuras tareas sobrehumanas; tercero, una critique de alto recall puede tolerar muchas sugerencias basura porque las personas las descartan, mientras que un assistant autónomo igual de poco confiable tendría un costo mucho mayor. El resultado es una proof of concept, no la demostración de que la scalable oversight esté resuelta.
\end{warningbox}

\teachervoice{El ponente describe el papel práctico de la critique como «ayudarte a hacer brainstorm de dónde deberías mirar». Esto es más moderado que «la IA ya conoce la respuesta correcta»: el modelo auxiliar quizá solo aporte direcciones de revisión, y el juicio final todavía requiere que la persona lo complete con base en el texto original, la gravedad y las especificaciones de la tarea.}

\subsection{Discriminator--critique gap: lo que el modelo sabe pero no dice}

En la sesión de Q\&A el problema avanza un paso más: el mismo pretrained model se ajusta por separado como discriminator y como critique model. El discriminator solo tiene que decidir qué response tiene una flaw; el critique model debe señalar la flaw en lenguaje natural. Si el primero es claramente más fuerte, eso indica que en las representaciones del modelo existe información útil para discriminar que no se le comunica de forma confiable a las personas.

Para convertir esta diferencia en una métrica comparable, puede definirse

\[
G_{\mathrm{D-C}}=\mathrm{Acc}_{\mathrm{disc}}-\mathrm{Help}_{\mathrm{critique}},
\]

donde $\mathrm{Acc}_{\mathrm{disc}}$ es la exactitud de discriminación del discriminator en la tarea de pares, $\mathrm{Help}_{\mathrm{critique}}$ es la tasa de éxito con la que la critique realmente ayuda a las personas a encontrar el defecto correcto, y $G_{\mathrm{D-C}}$ es el discriminator--critique gap. Un gap positivo grande indica que hay una distancia considerable entre «poder distinguir» y «poder o querer explicar».

\begin{importantbox}{Por qué el gap aporta más información que la calidad de la critique por sí sola}
Una critique deficiente puede deberse a que el modelo de plano no sabe, pero también a que sabe y no logra expresarlo con palabras, a que el objetivo de optimización no recompensa la divulgación completa o a que omite información de manera estratégica. El discriminator funciona como una sonda que aproxima una cota superior. Esa cota sigue limitada por la probe, los datos y la distribución, pero permite convertir la pregunta «cuánto sabe realmente el modelo» de un problema puramente filosófico en un experimento comparable.
\end{importantbox}

\subsection{Resumen de la sección}

Las targeted perturbations construyen una ground truth local mediante defectos artificiales; los critique experiments miden la ganancia real que aporta la assistance a la detección de errores por parte de las personas; el discriminator--critique gap examina la distancia entre las representaciones del modelo y el conocimiento que puede comunicar. Los tres son diagnósticos, no garantías definitivas.
\section{Q\&A de la clase: preferencias, seguridad, herramientas y problemas de alineación sin resolver}

La slide principal termina alrededor de 30:54, pero la mitad de la información de la clase proviene de la Q\&A posterior. Aquí no se transcribe mecánicamente en el orden en que aparecieron las preguntas, sino que la discusión se reorganiza en cinco temas de ingeniería: la división del trabajo entre humanos y máquinas, la uncertainty, el preference payload, la tool-assisted evaluation, y la outer/inner alignment junto con la interpretabilidad.

\subsection{Las Machines hacen el trabajo cognitivo; los Humans aportan la entrada de preferencias}

La visión de largo plazo del ponente no consiste en que las personas lean todo el material, escriban todas las pruebas y hagan todos los cálculos por sí mismas, sino en que la IA asuma ese cognitive labor escalable; las personas se concentran en la entrada que menos puede sustituirse: qué resultados vale la pena perseguir, qué tradeoff es aceptable y cuándo el riesgo es demasiado alto. Esta sección explica primero la división de responsabilidades de la última slide y después usa la Q\&A posterior para examinar en qué puntos esa división vuelve a trasladar la carga del juicio a las personas: la incertidumbre, la representatividad de las preferencias, los permisos de las herramientas y la confianza en la evidencia.

\lecturefigure{slide-17-machines-vs-humans.jpg}{La división del trabajo al final de la clase: las máquinas leen, calculan, escriben código y verifican; las personas comunican sus preferencias.}{Video oficial de Stanford Online 00:29:45--00:30:54.}

\begin{knowledgebox}{Lectura de la figura: la entrada de preferencias sigue requiriendo trabajo cognitivo}
``Humans: communicate preferences'' no significa solamente pulsar un botón de me gusta/no me gusta. Las personas deben aclarar objetivos, comparar valores en conflicto, decidir los estándares de evidencia, identificar riesgos inaceptables y juzgar si el assistant está manipulando el proceso de evaluación. El propósito de que las máquinas asuman más trabajo cognitivo es concentrar la atención humana, que es limitada, en estas decisiones que no pueden delegarse sin más.
\end{knowledgebox}

\subsection{Uncertainty calibration y la relación signal-to-noise de la critique}

Un estudiante pregunta primero si la hallucination puede resolverse mediante la incertidumbre del modelo. El ponente considera que enseñar al modelo a reconocer su incertidumbre cuando no está seguro ``debería poder resolverse'', pero en 2023 todavía no se había llegado a un estado satisfactorio. Muestrear varias veces tampoco es una respuesta completa: las muestras de un mismo fine-tuned model pueden estar muy correlacionadas, y preentrenar varios modelos independientes tiene un costo altísimo.

\begin{warningbox}{Qué es la calibration}
La calibration exige que, de un conjunto de eventos en los que el modelo declara un 80\% de certeza, aproximadamente el 80\% sea correcto; no equivale a agregar en el lenguaje un ``podría estar equivocado''. En los modelos generativos, además, hay que definir los eventos, descomponer las respuestas largas en claims verificables e impedir que el modelo eluda su responsabilidad con redacciones vagas. Expresar incertidumbre es útil, pero no sustituye la verificación de los hechos.
\end{warningbox}

La critique también tiene un problema de signal-to-noise. El ponente admite que muchas critiques son solo nitpicking o garbage; en un modo de alto recall, orientado a ``dar a las personas ideas para revisar'', esto puede ser aceptable, porque las personas descartan rápidamente las sugerencias irrelevantes. Pero si el assistant toma decisiones directamente, una precision baja consume atención, genera false alarms y hace que los problemas realmente graves queden sepultados en el ruido.

\subsection{Preference payload: las preferencias de quién, cómo se actualizan y cómo se evita el envenenamiento}

El RLHF suele representarse como ``human feedback'', pero la Q\&A indagó en un payload más difícil: a quién se considera human, cómo se representa a los distintos grupos y cómo se manejan las preferencias en conflicto. El ponente dijo abiertamente que, en ese momento, el conjunto de etiquetadores estaba formado en gran medida por las personas que la empresa podía contratar, y no era diverso ni representativo; algunas investigaciones convendría que las realizaran instituciones ajenas a las grandes empresas tecnológicas.

\begin{importantbox}{No entrenar en línea directamente con cualquier retroalimentación pública}
El ponente descartó primero un atajo peligroso: usar tal cual, para el entrenamiento continuo, cualquier entrada o retroalimentación de los usuarios en la interfaz del producto. Los ataques coordinados, las demostraciones maliciosas y el envenenamiento de datos convertirían el canal de retroalimentación en un plano de control. Un proceso seguro necesita, como mínimo, control de identidad y de calidad, detección de anomalías, muestreo estratificado, auditoría, conjuntos reservados y actualizaciones diferidas, en lugar de suponer que ``más retroalimentación es, por sí sola, más democrática''.
\end{importantbox}

Las preferencias también se desplazan con el tiempo. Volver a recolectar comparisons y hacer fine-tuning es relativamente barato en cómputo, pero la actualización del modelo cambia el estilo de sus salidas y su comportamiento, y las aplicaciones que dependen de prompts anteriores deben volver a adaptarse. Así, una alignment update es también un compatibility event: el sistema debe equilibrar un mejor ajuste a las nuevas preferencias con la estabilidad de los sistemas que dependen de él.

\begin{table}[H]
\centering
\caption{Cuatro límites de los preference data}
\begin{tabularx}{\textwidth}{L{2.5cm} X X}
\toprule
Límite & Pregunta central & Modo de falla \\
\midrule
Representation & ¿Quién entra al conjunto de etiquetadores? & Los grupos minoritarios y las necesidades especializadas se diluyen en el promedio \\
Aggregation & ¿Cómo se combinan las preferencias en conflicto? & Un conflicto normativo se registra erróneamente como ruido ordinario de etiquetas \\
Security & ¿Quién puede influir en el canal de retroalimentación? & Manipulación coordinada, envenenamiento y secuestro de la recompensa \\
Drift & ¿Cómo se actualizan las preferencias y el conocimiento con el tiempo? & Cambios bruscos de comportamiento, pérdida de compatibilidad de los prompts \\
\bottomrule
\end{tabularx}
\end{table}

\subsection{Browsing, API y external evidence}

Los estudiantes discutieron browsing, calculators, search y otras APIs. Las herramientas permiten que el modelo cite material externo, ejecute cálculos y verifique hechos en tiempo real, con lo que se pasa de ``responder con la memoria paramétrica'' a ``aportar evidencia verificable''. Pero las herramientas también amplían el action space: el modelo puede acceder a fuentes no confiables, filtrar datos, invocar operaciones con privilegios elevados o presentar un error de una herramienta como una conclusión segura.

\begin{knowledgebox}{Las herramientas extienden la alineación del texto a los límites del sistema}
Cuando el modelo solo genera tokens, el problema principal es la calidad de la respuesta; cuando puede invocar herramientas, también hay que definir credentials, dominios de red, permisos de escritura, presupuesto, registros, reversión y human approval. El tool result tampoco es ground truth por naturaleza: el orden de los resultados de búsqueda, el contenido de las páginas web, los datos de las API y el entorno de ejecución de código pueden manipularse. La supervisión escalable gana con herramientas más potentes, y también se vuelve más riesgosa por ellas.
\end{knowledgebox}

En el ámbito educativo, las herramientas también pueden provocar overreliance. Si los estudiantes se limitan a aceptar respuestas fluidas sin practicar el modelado, la deducción y la verificación, confundirán la capacidad de la herramienta con su propia comprensión. Un buen sistema debe exponer la evidencia, descomponer los pasos y animar al usuario a predecir y verificar, en lugar de entregar solo la conclusión final.

\subsection{Outer alignment, inner alignment y distribution shift}

En la parte final de la Q\&A se discute el inner optimizer. La outer alignment se ocupa de si la señal de entrenamiento externa expresa realmente los objetivos humanos; la inner alignment, de si el proceso de optimización que el modelo aprende internamente sigue persiguiendo el mismo objetivo fuera de la distribución de entrenamiento. La opinión personal del ponente es que hacer primero que la outer signal sea muy confiable podría ser especialmente importante, porque una señal confiable permite seguir entrenando y corrigiendo el optimizer interno en distribuciones nuevas.

\[
\underbrace{\text{trusted outer signal}}_{\text{what is rewarded}}
\quad+\quad
\underbrace{\text{coverage under shift}}_{\text{where it remains valid}}
\quad\Longrightarrow\quad
\text{continued correction of learned behavior}.
\]

Aquí, trusted outer signal designa la señal de evaluación en la que las personas están dispuestas a confiar; coverage under shift indica que esa señal sigue pudiendo juzgar la calidad en tareas y estados nuevos, y la flecha representa una esperanza de investigación, no una garantía demostrada.

\teachervoice{El ponente usa varias veces ``I think'', ``unclear'' e ``in practice''. Su argumento es que, si la señal externa es lo bastante confiable, parte del riesgo de inner alignment puede reformularse como un problema de distribution shift y de evaluación continua; no afirma que con ello la inner alignment haya desaparecido.}

\begin{warningbox}{La autocrítica no puede crear ground truth de la nada}
Hacer que el modelo revise sus propias respuestas, que dos modelos debatan entre sí o que un reward model evalúe la policy puede mejorar la cobertura; pero si todos los roles comparten el mismo punto ciego o el mismo objetivo equivocado, el ciclo solo se vuelve más coherente consigo mismo. Hay que conservar evidencia independiente, ejemplos adversarios, tareas verificables y failure signals que las personas puedan entender.
\end{warningbox}

\subsection{Interpretabilidad: útil, pero posiblemente neither sufficient nor necessary}

La discusión final fue muy mesurada. La interpretabilidad puede ayudar a detectar deception, a rastrear por qué un modelo da cierta respuesta y a formar parte del conjunto de herramientas de evaluación; pero aunque se logre explicar parte de los mecanismos internos, eso no indica automáticamente cómo convertir esas explicaciones en objetivos de entrenamiento confiables. Un riesgo más serio es el selection effect: si solo se descarta la misalignment que detectan las herramientas actuales, el entrenamiento o la selección pueden favorecer modos de falla más difíciles de detectar.

\begin{importantbox}{Los tres niveles del problema de la interpretabilidad}
El primer nivel es la detection: si es posible ver las características internas relacionadas con la falla; el segundo es la attribution: si es posible explicar cómo esas características producen el comportamiento; el tercero es el control: si es posible modificar el comportamiento de forma confiable sin dañar otras capacidades. Completar el primer nivel no equivale a completar el tercero; por eso la interpretabilidad puede ser muy útil y, aun así, no ser condición suficiente ni necesariamente la única vía necesaria.
\end{importantbox}

Por otro lado, si el comportamiento del sistema en todas las distribuciones relevantes puede evaluarse y restringirse de forma confiable, sigue abierta la pregunta de si sus ``razones'' internas deben ser completamente legibles. El ponente no dio una respuesta definitiva, sino que volvió a situar el objetivo en una ``really, really good evaluation signal'': poder juzgar de manera continua qué hizo realmente el modelo y, con base en ello, entrenarlo para que haga lo que de verdad queremos.

\subsection{Resumen de la sección}

La Q\&A amplía la alineación de una loss a un sistema de evaluación completo: hay que calibrar la incertidumbre, controlar el ruido de la critique, gobernar el preference payload, prevenir el envenenamiento, administrar los permisos de las herramientas, manejar el distribution shift y comprender los alcances de la interpretabilidad. Ninguna técnica aislada puede sustituir la trusted evaluation.
\section{Síntesis y ampliación}

La lógica de esta clase puede condensarse en una cadena cuya dificultad aumenta en cada eslabón. Primero, la human intent real no se limita a la instrucción literal; segundo, RLHF convierte parte de las preferencias implícitas en una señal optimizable mediante demonstrations, comparisons, un reward model y RL; tercero, InstructGPT y las primeras versiones de ChatGPT demostraron que esta señal mejora de forma notable la utilidad práctica, aunque persisten la hallucination, la prompt sensitivity y los problemas de representatividad; cuarto, cuando las tareas que completa el modelo rebasan la capacidad humana de evaluarlas de forma independiente, la señal del RLHF convencional se distorsiona; quinto, la AI-assisted evaluation busca que la máquina se encargue de leer, buscar, calcular y elaborar la critique, para que el ser humano se concentre en expresar sus preferencias; sexto, las targeted perturbations, los critique experiments y la discriminator--critique gap solo permiten medir paso a paso si esta línea de trabajo realmente amplía la supervisión.

\begin{importantbox}{Los cinco puntos más importantes de toda la clase}
\begin{enumerate}
  \item El objeto de la alignment es la human intent, no una obedience mecánica.
  \item El reward model es un proxy reutilizable de un conjunto limitado de datos de comparación y, al mismo tiempo, una nueva superficie de ataque expuesta a la optimización.
  \item La preference improvement, la capability, la truthfulness y la safety deben medirse por separado.
  \item Lo esencial de la scalable oversight no es que la IA decida los valores en lugar de las personas, sino transformar los objetos difíciles de evaluar en evidencia verificable.
  \item Todas las líneas de supervisión recursiva, self-critique, debate e interpretabilidad terminan enfrentando la misma pregunta: de dónde proviene una ground truth confiable o una trusted evaluation signal.
\end{enumerate}
\end{importantbox}

\subsection{Escalera de madurez de la evidencia}

El lector puede usar una escalera de cuatro niveles para juzgar la solidez de una alignment claim. El nivel de mecanismo explica cómo funciona el método; el nivel experimental actual indica si es eficaz en tareas controladas; el nivel de transferencia comprueba si se aplica a tareas más difíciles, más realistas y más adversariales; solo el nivel de garantía discute si sigue siendo confiable frente a sistemas más capaces que el evaluator. Los resultados de RLHF de esta clase se ubican principalmente en los dos primeros niveles; los resultados de scalable oversight son una proof of concept temprana del segundo nivel, y el papel final de la outer/inner alignment y de la interpretabilidad sigue siendo un tema de investigación abierto.

\begin{warningbox}{No presentar una línea de investigación como un hecho consumado}
La clase planteó una dirección prometedora: usar IA para ayudar a las personas a evaluar IA. No demostró que la curva de evaluación pueda alcanzar siempre a la curva de capacidad, ni que un mismo modelo vaya a revelar con honestidad todo lo que sabe. Unas notas de calidad deben conservar esta incertidumbre, porque la dificultad del problema radica justamente en que no podemos verificar directamente los sistemas potentes del futuro con los experimentos pequeños de hoy.
\end{warningbox}

\subsection{Lecturas adicionales}

Todos los materiales siguientes se publicaron antes de la fecha de la clase y respaldan directamente su línea principal:

\begin{itemize}
  \item \href{https://arxiv.org/abs/1706.03741}{Christiano et al., Deep Reinforcement Learning from Human Preferences}: preferencias por pares de trayectorias y reward learning en línea.
  \item \href{https://arxiv.org/abs/1811.07871}{Leike et al., Scalable Agent Alignment via Reward Modeling}: agenda de investigación de reward modeling y recursive assistance.
  \item \href{https://arxiv.org/abs/2203.02155}{Ouyang et al., Training Language Models to Follow Instructions with Human Feedback}: implementación pública y evaluación de SFT, RM, PPO y PPO-ptx.
  \item \href{https://arxiv.org/abs/2009.01325}{Stiennon et al., Learning to Summarize with Human Feedback}: optimización de preferencias en la generación de lenguaje.
  \item \href{https://arxiv.org/abs/2206.05802}{Saunders et al., Self-critiquing Models for Assisting Human Evaluators}: critique assistance, targeted flaws y discriminator--critique gap.
  \item \href{https://arxiv.org/abs/1805.00899}{Irving et al., AI Safety via Debate}: otra línea que amplifica el juicio humano mediante el diálogo y la argumentación adversarial.
\end{itemize}

En última instancia, la alignment no consiste en conectar un «goodness score» a la salida del modelo, sino en diseñar una cadena de retroalimentación que no pierda confiabilidad a medida que crece la capacidad. Al final, el ponente no prometió que algún componente resolvería todos los problemas; la postura de investigación más confiable que propuso es: hacer medibles las tareas de evaluación, hacer verificable la evidencia, hacer explícitos los límites de falla, y seguir preguntando qué sabe el modelo, qué dice y qué hace en realidad.

\subsection{Autoevaluación posterior a la clase}

Después de leer esta clase, no basta con saber repetir el flujo de RLHF; también hay que poder determinar qué tramo de la cadena de retroalimentación mejora realmente una nueva propuesta de alignment. Las cuatro preguntas siguientes sirven como lista mínima de verificación al leer artículos, diseñar experimentos o revisar las afirmaciones de un producto.

\begin{knowledgebox}{Cuatro preguntas para examinar una nueva alignment claim}
\begin{enumerate}
  \item ¿La señal de entrenamiento proviene de demonstrations, comparisons, verifiable outcomes o del juicio de otro modelo no validado?
  \item ¿El experimento mide la preference, la tasa real de aciertos, la calibration, la robustness, o solo un proxy fácil de optimizar?
  \item Cuando la dificultad de la tarea rebasa la capacidad humana de evaluación independiente, ¿qué evidence verificable aporta el sistema, más allá de una explicación más fluida?
  \item Cuando ocurre una falla, ¿es posible distinguir entre que el modelo no sabe, que sabe pero no puede expresarlo, que sabe pero no quiere decirlo, y que el evaluator no lo detectó?
\end{enumerate}
\end{knowledgebox}

\end{document}
